\chapter{Discussion}
\label{ch:discussion}

The results of the previous chapter are, on their surface, a collection of specific findings: a step size that does not transfer between models, an annealing artifact that mimics convergence, a Metropolis--Hastings correction that rescues one sampler and destroys another, a gradient that carries no information, a likelihood trap in which the best-scoring text is the worst text, three distinct modes of GFlowNet collapse, and a constraint that does not steer. Taken one at a time, each is a result about a method. Taken together, they are evidence for a single claim about the object all of these methods share. This chapter draws them into that single account. It first states explicitly which research question each part of the results answers and in which section, so that the mapping from questions to evidence is unambiguous. It then presents the unified explanation and argues for its strength on the grounds that it is reached by two independent routes. It connects each finding back to the specific prior work of Chapter~\ref{ch:related}, so that the contribution can be located precisely against the literature. It states carefully what the results do and do not license one to conclude. And it closes with the experiment that would falsify the central explanation, which is, in the author's view, the most important thing the thesis leaves for future work.

\section{Answers to the Research Questions}
\label{sec:disc-rqs}

It is useful to state plainly which research question each result answers and in which section of Chapter~\ref{ch:results} the answer is established, both because the reader deserves an explicit map and because the discipline of writing the map down is a check that every question posed in the introduction has in fact been answered by evidence rather than by assertion.

\textbf{RQ1} asked whether the frozen likelihood can be sampled effectively with faithful Langevin dynamics on the discrete manifold, and if not, why. It is answered in the negative across three sections. Section~\ref{sec:results-stepsize} shows that the sampler cannot even be made to move without model-specific calibration, which already denies the strong plug-and-play reading. Section~\ref{sec:results-fallacy} shows, through the direction-versus-magnitude ablation, that once the sampler does move, following the true gradient is statistically indistinguishable from following a random direction of equal norm, so the gradient is not a usable search direction. And Section~\ref{sec:results-linradius} supplies the mechanism behind that fact, namely that the linearization radius of the surrogate the sampler uses is smaller than the smallest admissible discrete step, so the gradient is inapplicable rather than merely imprecise, and is compounded by the structural blindness of the input-embedding gradient to the self-term of the energy change. The three sections together answer both the ``whether'' and the ``why'' of RQ1.

\textbf{RQ2} asked about the role of the Metropolis--Hastings correction in each sampler. It is answered in Sections~\ref{sec:results-quench} and~\ref{sec:results-mh}, and the answer is that the correction plays opposite roles in the two samplers, which is itself informative. For the discrete sampler the correction is beneficial: Section~\ref{sec:results-quench} shows that it removes the quenching artifact and produces smooth early convergence, so here theoretical correctness and empirical performance align. For the continuous sampler the same correction is catastrophic: Section~\ref{sec:results-mh} shows that it rejects almost every useful move, and the decomposition of the acceptance ratio attributes the collapse specifically to the proposal term, connecting it to the non-Lipschitz drift induced by nearest-neighbour projection. The correction does not make the continuous sampler work; it reveals that the continuous sampler was only ever functioning as a biased optimizer, and that its apparent successes were successes of optimization under early stopping rather than of sampling.

\textbf{RQ3} asked whether amortization with a GFlowNet escapes the failure. It is answered in Section~\ref{sec:results-gfn}, and the answer is no. Three distinct failure modes appear, capacity collapse, length collapse, and reward hacking, and the unifying experiment shows that Langevin sampling on the amortized energy is indistinguishable from sampling on the untuned base, which locates the cause in the energy rather than in the search procedure.

\textbf{RQ4} asked whether an additive constraint can steer generation on this landscape. It is answered in Section~\ref{sec:results-constrained}, and the answer is no: the steering gains are small and inconsistent in sign, and the failure follows from the same defect established for the fluency energy rather than being a separate phenomenon.

\section{A Single Defect, Reached by Two Routes}
\label{sec:disc-unified}

The organizing claim of this thesis is that every failure catalogued above is a surface manifestation of one underlying defect: the frozen likelihood of an autoregressive language model is not a usable energy function for sampling, because the training objective that produces it never shapes it to be one. The argument for this claim was made incrementally through the results; here it is stated as a whole.

Recall the distinction drawn in Section~\ref{sec:bg-ebm}. Teacher forcing optimizes each local conditional $p_{\text{LM}}(x_t \mid x_{<t})$ to predict the next token given a correct prefix drawn from the data. That is the entire content of the objective, and modern models satisfy it superbly. What the objective never does is constrain the behaviour of the joint log-likelihood $\log p_{\text{LM}}(\bm{x})$ as a function of $\bm{x}$ when $\bm{x}$ is varied off the data manifold, which is exactly the regime in which every energy-based sampler operates, because sampling is precisely the business of moving through sequence space, including through its off-distribution regions, in search of good sequences. The joint likelihood is therefore never asked to be smooth, never asked to be a calibrated measure of quality, and never asked to have a gradient that points toward better sequences. Energy-based decoding borrows this quantity and treats it as though it had all three properties. The results show, in three complementary ways, that it has none of them: it is not smooth, as the discontinuous Voronoi landscape of Section~\ref{sec:results-mh} demonstrates; it is not a good measure of quality, as the likelihood trap of Section~\ref{sec:results-trap} demonstrates; and its gradient does not point toward better sequences, as the gradient fallacy of Section~\ref{sec:results-fallacy} and the linearization failure of Section~\ref{sec:results-linradius} demonstrate. These are not three separate deficiencies but three faces of one absence, the absence of any training pressure toward a well-behaved global energy.

There is a cleaner and more mechanical way to see why the \emph{gradient} in particular is the defective part, and it follows directly from the decomposition of Section~\ref{sec:results-linradius}. The only term in the likelihood that registers how well a token fits its own left context enters through a discrete index into the output softmax, not through the continuous input embedding that the sampler differentiates. The input-embedding gradient can therefore see a substituted token only through that token's effect on the \emph{later} tokens it feeds forward, and it is structurally blind to the token's own fit. This blindness is partial in the interior of a sequence, where a token still shapes its successors, and it becomes total at the final position, where the token's embedding feeds no realized prediction at all: the gradient of the sequence log-likelihood with respect to the last token's input embedding is exactly zero, even though the energy difference between candidate final tokens is precisely the model's conditional log-ratio and is therefore maximally informative. At that position the energy knows exactly which token is best and the gradient provably cannot, which is the sharpest possible statement of the whole diagnosis: what fails is not the information in the energy but the attempt to extract a search direction from its input-embedding gradient. Making this last-position argument exact, and pairing it with an energy-only sampler that succeeds on the same task where the gradient arm fails, is the natural culmination of the diagnosis and the immediate next analysis the work points to.

The strength of the evidence for this claim rests on a feature of the study's design that is worth making explicit, because it is what elevates the conclusion above a collection of negative results. The claim is reached by two routes that share no machinery. The gradient-guided route differentiates the likelihood, and it fails: the gradient is uninformative because its surrogate is uncorrelated with the true energy change, and the correction, applied honestly, rejects every useful move. The amortized route never touches the gradient at all; it treats the likelihood only as a scalar reward, and it fails too, because the reward has no well-behaved optimum where coherent text lies, so the learned policy collapses in whichever direction the reward's implicit incentives push it. Two methods with nothing in common beyond the energy they use should not fail in ways that trace back to the same property of that energy, unless the energy is indeed the common cause. And the unifying experiment removes even the last doubt: running the gradient-guided sampler on the amortized energy reproduces exactly the behaviour of running it on the base, so passing the energy through the amortization changes nothing about the sampler's fate. The energy, and not the sampler, and not the search procedure, is the problem.

The likelihood trap ties the two routes together at the level of intuition, and it explains an otherwise puzzling feature of the whole picture. Because low energy is degenerate text, the goal that both methods pursue is the wrong goal, and this casts the failure of the gradient-guided samplers in an unexpected light. Their inability to reach low energy, which the earlier sections established as a failure, is in a sense a reprieve, because reaching low energy, as both the quenching phase of the discrete sampler and the reward hacking of the GFlowNet show, produces bad text. The methods are protected from their own objective by their inability to optimize it, and the moments when they do optimize it effectively, beam search on the length-normalized GFlowNet being the clearest, are exactly the moments when their output degenerates into strings like the ones displayed in Section~\ref{sec:results-trap}. A method that worked, in the narrow sense of reaching the energy's minimum, would produce gibberish; a method that produces fluent text does so by failing to reach the minimum. This is the deepest sense in which the frozen likelihood is not an energy worth sampling.

\section{Connections to the Related Work}
\label{sec:disc-related}

The findings of this thesis stand in a definite relationship to the literature reviewed in Chapter~\ref{ch:related}, and making that relationship explicit locates the contribution precisely rather than leaving the reader to infer it.

The likelihood trap measured in Section~\ref{sec:results-trap} is a direct, quantitative reproduction on the study's own models of the neural-text-degeneration phenomenon that \citet{holtzman2020curious} identified: that the most probable sequences under a neural language model are degenerate, and that maximization-based decoding produces worse text than sampling. Chapter~\ref{ch:related} presented that as a known hazard of decoding. The contribution here is to re-derive it as a structural obstacle for energy-based sampling specifically, because in the energy-based framing the degeneracy of high-probability text is not merely a decoding pitfall to be avoided but a statement that the \emph{target} of the sampling procedure is undesirable, which is a more damaging thing to say about a method that is defined by its pursuit of that target. The monotone likelihood-repetition relationship measured here is the quantitative form of that statement.

The anisotropy results of Section~\ref{sec:results-aniso} confirm, on GPT-2 Large and Llama-3, the representation-degeneration and embedding-anisotropy phenomena documented by \citet{gao2019representation} and \citet{ethayarajh2019contextual}. Chapter~\ref{ch:related} raised anisotropy as a reason to distrust Euclidean distance in embedding space, and the measurements here quantify that distrust. But they also draw a consequence those works did not consider: because the anisotropy scale differs by roughly a factor of three between the two models, a single step-size schedule cannot transfer between them, which is the concrete and previously unremarked cause of the calibration failure that opened the results chapter. The anisotropy literature explains why a token can be close in embedding space yet wrong; this thesis adds that the same geometry makes the samplers model-specific in a way that undercuts their claim to be plug-and-play.

The Metropolis--Hastings breakdown of Section~\ref{sec:results-mh} is the empirical counterpart to a gap noted in Section~\ref{sec:rel-energy}: that the energy-based decoding literature frequently omits the correction, and that where it is included, the difficulty of computing the reverse proposal on a projected landscape is seldom examined. This thesis takes the correction seriously, implements it faithfully following the convergence theory of \citet{roberts1996exponential}, and measures exactly what the common omission was concealing. Seen through this lens, methods such as COLD \citep{qin2022cold} and MuCoLa \citep{kumar2022gradient}, which the continuous sampler of this thesis follows, are best understood not as samplers from the stated energy but as biased optimizers, and their reported successes are successes of optimization under early stopping rather than of sampling. This is not a criticism of those works, which are careful about what they claim, but a clarification of what the methods actually do, and it is a clarification the correction makes unavoidable. The same lens sharpens the relationship to the annealed stochastic-gradient Langevin dynamics of \citet{welling2011bayesian}, whose guarantee is that annealing the step size toward zero makes the chain converge to the target distribution. The quenching effect of Section~\ref{sec:results-quench} is what that recipe degenerates into once the correction is dropped: the annealed noise does not deliver a fair draw from the distribution but freezes the chain into whichever mode it last occupied, so the very schedule that the theory of \citet{welling2011bayesian} relies on becomes, without the correction, the mechanism of a biased optimization.

Finally, the GFlowNet results of Section~\ref{sec:results-gfn} qualify, without contradicting, the picture presented by \citet{hu2024amortizing}, whose formulation and codebase this thesis builds on. Their work showed that amortized sampling from a frozen-model posterior can yield diverse, transferable samples on tasks where the reward is well-behaved, and nothing in this thesis disputes that. The narrower and more critical finding here is that when the reward is the raw frozen likelihood used as an energy over free-form sequences, the amortization inherits the pathology of that energy and the learned policy collapses. The two findings are compatible, and the axis that reconciles them is precisely the one this thesis foregrounds: whether the reward defines a landscape with a coherent-text optimum. Where it does, amortization is useful; where it is the raw autoregressive likelihood, it is not, and the contribution here is to show, by the unifying experiment, exactly which case the frozen likelihood falls into.

Finally, the gradient-free samplers isolated in Section~\ref{sec:rel-energy}, Mix-and-Match \citep{mireshghallah2022mixmatch} and the twisted sequential Monte Carlo of \citet{lew2023sequential} and \citet{zhao2024probabilistic}, stand in a confirming rather than a contradicting relationship to the central result. Each samples from a frozen-model energy without differentiating it, and each works. Chapter~\ref{ch:related} raised this as a hint; the Gibbs baseline of Section~\ref{sec:results-baselines} converts the hint into direct evidence by running a gradient-free sampler on the identical energy the Langevin methods use and recovering coherent text where they cannot. Taken together, the external methods and the in-platform Gibbs result draw the same line: the frozen likelihood is a serviceable object to \emph{evaluate}, and to search without gradients, and an unusable object to \emph{differentiate}. This is the empirical warrant for the narrowed claim of the conclusion, and it is why the thesis does not overreach into asserting that the energy is useless, only that its gradient is.

\section{What the Results Do and Do Not Show}
\label{sec:disc-scope}

It is important to state the limits of these conclusions precisely, both out of scientific honesty and because a diagnosis is only useful if its scope is clear. The thesis shows that the frozen likelihood of a \emph{standard autoregressive} language model is a poor energy for gradient-guided and amortized sampling on discrete text, and it explains why, in terms of the training objective and the geometry of the token embedding space. It does not show that energy-based controllable generation is impossible in general. It shows that a particular and widely used instantiation of it fails, and it identifies the specific property responsible for the failure, which is a more useful result than a blanket impossibility claim precisely because it points to what would have to change.

Several caveats bound the empirical claims, and they are stated here rather than buried. The KL-divergence metric depends on the position of the corrupted token within the sequence, which is why, as Section~\ref{sec:meth-metrics} explained, the central results are argued from the absence of a gap between conditions rather than from small measured differences; a null gap is robust to a positionally noisy metric in a way that a narrow measured win would not be, and the reader should read the gradient-fallacy result as ``no reliable difference'' rather than as a precise numerical equality. The Llama-3 results support only qualitative cross-architecture claims, because the tokenizer and embedding scale differ from the GPT-2 family, so the Llama measurements establish that the phenomena occur across architectures without licensing a numerical comparison between the two. The linearization result, as Section~\ref{sec:results-linradius} was careful to note, is that the surrogate is uniformly uninformative across the admissible range of distances rather than informative up to a sharp radius and useless beyond it; the language of a ``linearization radius'' is a convenient summary of a decay, not a claim of a clean threshold, and the honest reading of Figure~\ref{fig:lin-radius} is that the correlation is low everywhere the sampler operates. And the constrained-generation extension establishes the absence of reliable steering on this energy, not the impossibility of steering on a better one.

The thesis is also candid about the relationship between its title and its content, because that candour is part of the diagnostic honesty the work aims for. The original objective was controllable generation via faithful Langevin dynamics, conceived as a two-stage plan in which the first stage would establish that energy-guided sampling on a frozen model is viable and the second would add the constraint term and perform sentiment and toxicity control. The first stage did not hold. Faced with that, there were two paths: to proceed to the second stage regardless, building control on a foundation the study had shown to be unsound and thereby producing numbers whose meaning could not be interpreted, or to turn to the more valuable question of why the foundation fails. The work took the second path. What it delivers is therefore a systematic diagnosis of the failure of a prerequisite, with the mechanisms identified, measured, and cross-checked, and with the constrained-generation extension included as direct evidence of the downstream consequence, rather than a demonstration of successful control. A diagnosis of why a widely pursued approach does not work, when it is supported by measurement and reaches a single explanation by independent routes, is in the author's view a more useful contribution to the field than a marginal and uninterpretable positive result would have been.

\section{Implications and the Falsifying Experiment}
\label{sec:disc-future}

The most valuable consequence of locating the defect in the training objective, rather than in the sampler or the search procedure, is that it makes a sharp and testable prediction, and a diagnosis that makes a falsifiable prediction is worth more than one that does not. If the reason the frozen autoregressive gradient is useless is that teacher forcing never shapes a score function over sequences, then a model whose training objective \emph{does} shape such a function should not suffer the same fate, and running the identical protocol of this thesis on such a model would either confirm the diagnosis or refute it.

\textbf{Diffusion language models} are exactly such models \citep{li2022diffusion, lou2024discrete}. A diffusion model is trained, by construction, to denoise corrupted inputs, and denoising is equivalent to learning the gradient of the log-density of the data, the score function, by the identity of \citet{vincent2011connection} that also underpins score-based generative modeling \citep{song2019generative, ho2020denoising}. Where the autoregressive model is trained only to predict the next token and is never asked about the shape of its energy off the data manifold, the diffusion model is trained precisely to know, at every noise level, in which direction the data density increases, which is exactly the information the Langevin sampler needs and exactly the information the autoregressive gradient was found to lack. Running the diagnostic protocol of this thesis on a diffusion language model therefore constitutes a positive control in the strict experimental sense. If the linearization correlation is high where it was zero, if the gradient outperforms a random direction where it did not, and if the sampler recovers coherent text where it produced degeneracy, then the training objective is isolated as the causal variable and the diagnosis of this thesis is confirmed. If, on the other hand, the diffusion model fails in the same ways, then the diagnosis is wrong, the cause lies deeper than the training objective, and a different explanation must be sought. Either outcome is informative, which is the mark of a well-posed control, and the experiment is inexpensive because score-based discrete diffusion models that share the GPT-2 tokenizer are available off the shelf. This is, in the author's assessment, the single most important next step, and the design of this thesis, in particular its reliance on a masked-recovery task that a diffusion model can perform without modification, was shaped throughout to make this test clean and immediate.

The concrete instrument for the control is the score-entropy discrete diffusion model of \citet{lou2024discrete}, which shares the GPT-2 tokenizer and can therefore be placed inside the existing masked-recovery harness without retokenization. It is honest to report the status of this experiment precisely, because it is designed but not yet run. The measurement pipeline for it has been built and validated end to end against a stand-in model, so that the linearization diagnostic of Section~\ref{sec:results-linradius}, computing the surrogate-to-truth correlation on the diffusion model's learned score in place of the autoregressive gradient, executes and emits its outputs correctly; the two model-specific hooks that read the diffusion score are marked in the code for verification against the reference implementation. What remains is the single substantive step of loading a trained checkpoint and running it, which is deferred here rather than reported, so that no claim in this thesis rests on a result not yet in hand. The prediction the control would test is sharp. If the diagnosis is right, the linearization correlation that was negligible for the autoregressive gradient should be substantially positive for the diffusion score, the true gradient should outperform a random direction of equal norm where before it did not, and the sampler should recover coherent text where it produced degeneracy. If instead the diffusion model fails in the same ways, the training objective is exonerated and the cause must lie deeper than this thesis locates it. Because both outcomes would be informative, the experiment is a genuine control and not a confirmation-seeking exercise, and running it is the most valuable single step this thesis leaves open.

A second and complementary direction is to attack the defect at its source rather than route around it. If the problem is that the likelihood is not a score function, then one option is to make it one, by fine-tuning the model so that its sequence log-density is aligned directly with a target energy using a score-matching objective. This is a lighter intervention than the reinforcement-learning-style training of a GFlowNet, and it is directed precisely at the identified deficiency rather than at a downstream symptom of it. Only once an energy landscape is known to be navigable does the additive constraint term of the plug-and-play framework become meaningful, because only then is there a fluency energy for the constraint to be added to that a sampler can actually follow, and at that point the original goal of controllable generation becomes reachable by the route this thesis set out to take. The contribution of the present work to that longer programme is a map of where the ground gives way, drawn carefully enough that the next traveller knows both which step is unsafe and why, together with a falsifiable hypothesis about what firmer ground would look like.
